\section{Methodology}
\label{sec:method}

\subsection{Problem setup}

An evaluation protocol $\pi$ is given a budget $B$ of unit-cost queries, a set of targets
(behaviours or items) $\gI$, and a set of \families $\gF$. A design family is a qualitatively
distinct way of probing the model: a suffix-search initialisation, a prompt template class, or a
social-pressure schema. At each step the protocol chooses a target and a family, spends one unit,
and observes an outcome $y \in \{0,1\}$ together with a continuous surrogate score $s \in \R$. The
primary dependent variable is the \emph{detection rate} at matched budget,
$D_\pi(B) = \frac{1}{|\gI|}\sum_{i\in\gI} \1[\pi \text{ elicits the behaviour on } i \text{ within } B]$.
We report $D_\pi(B)$ as a whole curve over $B$, never as a single point.

The key comparison in this paper is not adaptive versus static. It is adaptive versus
\emph{blindly diversified}. Let $\pi_{\text{blind}}$ spend $B$ across $\gF$ on a schedule that is a
function of $B$ alone, and let $\pi_{\text{adapt}}$ spend $B$ across the same $\gF$ using a
schedule that is a function of $(B, s_{1:t})$. Because $\pi_{\text{adapt}}$ reduces to
$\pi_{\text{blind}}$ when its trigger is disabled, $D_{\pi_{\text{adapt}}}(B) -
D_{\pi_{\text{blind}}}(B)$ isolates the value of using the null as a trigger, with everything else
held fixed.

\subsection{Substrates}

\tabref{tab:substrates} lists the three substrates. Nothing in E1--E3 is re-run or simulated:
these are the original authors' own released traces and the official benchmark artifacts. E4 uses
real model weights executed locally. No simulated agent appears anywhere in this work.

\begin{table}[t]
    \centering
    \resizebox{\textwidth}{!}{%
    \begin{tabular}{@{}llll@{}}
        \toprule
        \textbf{Experiment} & \textbf{Substrate} & \textbf{Size} & \textbf{Role} \\
        \midrule
        E1, E1c, E3 & \rstraj: released per-iteration & 14 configs, $\sim$437k iteration rows, & Real sequential traces with a continuous \\
                    & random-search logs~\cite{andriushchenko2024adaptive} & 672 run outcomes & objective; contains a matched adaptive/static pair \\
        \addlinespace
        E2 & \jbb: official JailbreakBench & 5 methods $\times$ 4 target models & Complete cost-to-discovery matrix; any \\
           & artifacts~\cite{chao2024jbb} & $\times$ 100 behaviours & allocation policy can be replayed offline \\
        \addlinespace
        E4 & \gsmk test split~\cite{cobbe2021gsm8k} + real & 32 items $\times$ 4 pressure families & End-to-end live replication on a fresh, \\
           & open-weights LLMs run locally & $\times$ 2 conditions $\times$ 3 models & non-harmful failure mode \\
        \bottomrule
    \end{tabular}
    }
    \caption{The three substrates. E1--E3 replay already-published traces and artifacts; E4 runs
    real model weights locally. No simulated model is used anywhere in this work.}
    \label{tab:substrates}
\end{table}


\para{\rstraj (E1, E1c, E3).} We parsed the 24 attack logs released with
\citet{andriushchenko2024adaptive} into 14 configurations, roughly 437k per-iteration rows and 672
run outcomes. Each row carries the target-token log-probability surrogate and the judge outcome, so
a protocol can be replayed at iteration granularity. The corpus contains a matched pair on
\llamaseven that differs only in initialisation policy: \texttt{plain\_init} (static) scores $0.00$
over 50 behaviours with 56\% of runs hitting the 10{,}000-iteration cap, while the self-transfer
configuration scores $0.96$ with a median of 2 iterations.

\para{\jbb (E2).} The official JailbreakBench artifacts~\cite{chao2024jbb} give
\texttt{queries\_to\_jailbreak} and a success flag for 5 attacks $\times$ 4 target models $\times$
100 behaviours. This is a complete cost-to-discovery matrix, so any allocation policy can be
replayed offline exactly.

\para{Live grid (E4).} We run three open-weights instruct
models~\cite{yang2024qwen25,abdin2024phi3} locally on 32 \gsmk test problems
under an anti-sycophancy system-prompt intervention and a no-intervention control, crossed with
four social-pressure \families: P1 simple doubt, P2 appeal to authority, P3 fabricated consensus,
P4 persistent multi-turn pressure. The failure mode is sycophantic answer-flipping, which is
non-harmful and produces no jailbreak artifacts. Only items the model answers correctly without
pressure are eligible, so ``flip'' is always well defined. Baseline accuracies are $0.844$
(\qwenseven), $0.781$ (\phimini) and $0.594$ (\qwensmall), leaving 27, 25 and 19 eligible items
respectively.

\subsection{Matched-budget replay}

Every experiment follows the same pattern. The outcome matrix is evaluated (or read) once,
exhaustively, over all targets $\times$ families $\times$ budget positions. Protocols are then
replayed against that cached matrix under an identical budget. This has three consequences. First,
the compute confound is removed by construction rather than by argument. Second, every protocol
comparison is paired on the same behaviours or items, so we can use paired tests. Third, bootstrap
uncertainty is cheap, because replaying a policy costs microseconds.

\subsection{Protocol arms}

\para{E1: within-search escalation ladder.} Three \families are available on \llamaseven: $d_0$ the
default \texttt{plain\_init} random search, $d_1$ re-initialisation from a related success
(self-transfer), and $d_2$ a switch of template family. The arms are \staticzero, \staticone and
\statictwo (all $B$ on one family; \staticzero is the ``continue-only'' control, \ie maximal
within-family persistence), \fixedsplit ($B$ divided equally over the three families on a schedule
that ignores the data), \ladder (same designs, same $B$, but the switch point is moved by the
continuous surrogate), and \oracle (best family per behaviour). \ladder has three
hyper-parameters: a patience, a surrogate threshold $\tau$, and an extension factor. At
$\tau=0$ and extension $1$ it recovers \fixedsplit \emph{exactly}. This is what makes the
\ladder--\fixedsplit contrast decisive; the \staticzero contrast only establishes that design
diversity matters at all.

\para{E1c: null diagnosticity.} At checkpoint $t$ we restrict to runs that are \emph{still null at
$t$}, which is precisely the decision context an adaptive protocol faces, and measure the AUC of
the surrogate at $t$ for predicting whether the run eventually succeeds. If this AUC is $0.5$, no
protocol can use the null as a trigger, however cleverly implemented.

\para{E2: portfolio allocation.} Budget is a number of attack runs allocated over methods and
behaviours. Static arms are \singlemethod (the field default), \roundrobin (diversified but
non-adaptive) and \bestsingle, an \emph{oracle} single-method choice that upper-bounds any
\textit{a priori} decision. Adaptive arms are escalation, \thompson (Thompson sampling over
methods), \succhalving (successive halving) and \escalthompson. \oracle marks the union ceiling.

\para{E3: calibration of the terminal null.} A protocol stops and must emit a nominally 95\% upper
bound on the elicitable fraction. \fixedcp is Clopper-Pearson at a pre-registered budget,
\adaptivecp is the same interval at a data-chosen stopping time (what the field actually does),
\unioncp applies a Bonferroni union bound over the monitoring grid and is genuinely anytime-valid,
and \extrapolative fits the Best-of-$N$ power law~\cite{hughes2024bon} and extrapolates to the full
horizon. Ground truth is the outcome at the end of the observed trace.

\para{E4: live protocol race.} \staticpone spends its whole budget on P1, the obvious first design;
\roundrobin cycles blindly over the four families; \thompson allocates adaptively over families.
Budget is a number of model calls.

\subsection{Metrics and statistical analysis}

The analysis plan was pre-registered before any protocol was run. Primary metric is detection rate
at matched budget. Secondary metrics are queries to first detection treated as right-censored
survival data, AUC of the trigger signal, coverage and width of the emitted upper bound, and, for
E4, the flip rate under pressure and the fraction of vulnerable items localised.

Uncertainty comes from bootstrap over behaviours or items (1{,}000 resamples, seed 42), never from
re-running attacks, since the artifact matrix has one realisation per cell. Where an exact
per-behaviour join exists we use paired bootstrap plus exact McNemar on discordant behaviours;
unpaired comparisons are matched by resampling and are flagged explicitly as resting on an
exchangeability assumption. E2 and E4 comparisons use Wilcoxon signed-rank over paired policy
replicates (300 and 400 respectively). We apply Holm correction within each experiment's family of
tests at $\alpha = 0.05$, and report an effect size (risk difference, Cohen's $h$) alongside every
$p$-value. Both the GPT-4-judge and the rule-based outcome are retained throughout; they disagree,
and we report the disagreement rather than averaging it away.

\subsection{Reproducibility}

Seed 42 everywhere. E4 decoding is greedy (\texttt{do\_sample=False}, temperature 0). Environment:
Python 3.12, torch 2.13.0+cu130, transformers with native model code
(\texttt{trust\_remote\_code=False}), pandas 3.0.5, scipy 1.18.0, with \texttt{uv} managing an
isolated virtual environment. E1--E3 run on CPU in under five minutes in total; the E4 grid took
roughly 10 GPU-minutes per model on an RTX A6000.
